\documentclass[11pt]{article}
\usepackage[T1]{fontenc}
\usepackage{lmodern,amsmath,amssymb,amsthm,microtype,xcolor}
\usepackage[letterpaper,margin=1in]{geometry}
\usepackage[colorlinks=true,linkcolor=blue!45!black,urlcolor=blue!45!black]{hyperref}
\newtheorem{theorem}{Theorem}[section]
\newtheorem{lemma}[theorem]{Lemma}
\newtheorem{proposition}[theorem]{Proposition}
\newtheorem{corollary}[theorem]{Corollary}
\newtheorem{remark}[theorem]{Remark}
\setlength{\emergencystretch}{3em}
\title{A strict cycle-count deficit at bounded girth}
\author{Collaborative research working draft}
\date{11 November 2026}
\begin{document}
\maketitle
\begin{abstract}
Fix $g\ge5$.
Every connected loopless cubic multigraph of rank $r$, girth $g$, and no cyclic cut of size at most four satisfies
\[
 c(G)\le2^{r-1}-2^{r-g-1}+o_g(2^r).
\]
The proof partitions terminal paths by their first departure from a shortest cycle.
Its complement may be disconnected.
Every exterior component contains a cycle, which gives the required rank decrease after each prefix deletion.
Consequently, every remaining above-half sequence without small cyclic cuts has girth tending to infinity.
The argument also settles bounded-rank shores for every fixed cut size.
The unrestricted $O(r2^{r/2})$ error remains open.
\end{abstract}

\section*{Status and inputs}
The proof has passed four independent full mathematical reviews within the research team.
It is not formalized in Lean, and publication priority has not been established.
The geometric prefix argument and first-divergence count are proved here.
The preceding research packets supply the following analytic inputs.

Rank means $|E|-|V|+\operatorname{comp}$.
A normalized $k$-pole has $k$ distinct degree-two ports and all other degrees three.
A cyclic cut has a cycle on each shore.
Parallel edges remain distinct, and a parallel pair is a cycle.
Write $\psi(r)$ for the maximum cycle count at rank $r$.

\section{Established analytic inputs}
\begin{theorem}[Asymptotic cycle bound]
Uniformly over connected graphs of rank $r$,
\[
 c(G)\le(1/2+o(1))2^r.
\]
The stronger unrestricted estimate proved in the accompanying packet is
\[
 \psi(r)\le2^{r-1}+O\!\left(2^r(\log_2 r)^{11/2}/\sqrt r\right).
\]
\end{theorem}

\begin{theorem}[High-density paths force a small cyclic cut]\label{thm:high-density-pair-small-cyclic-cut}
For $0<\varepsilon\le1/2$, let $C_\varepsilon$ be the established additive terminal-path stability constant.
Let $X$ be a normalized $k$-pole of rank at least
\[
 (k+1)\bigl(9\lceil C_\varepsilon\rceil+8k+9\bigr).
\]
If some terminal pair has at least $(1/2+\varepsilon)2^{\beta(X)}$ paths, then $X$ contains two disjoint cyclic vertex sets.
They avoid all ports and have boundary at most four in every cubic completion.
\end{theorem}

\begin{theorem}[Protected version]\label{thm:high-density-protected-cyclic-sets}
For any fixed number $m$ of protected vertices, the preceding conclusion can also avoid those vertices.
A sufficient rank threshold is
\[
 (k+m+1)\bigl(9\lceil C_\varepsilon\rceil+8k+4m+9\bigr).
\]
\end{theorem}

The path-stability proof appears in \nolinkurl{terminal-path-stability.pdf}.
The small-cut and protected refinements appear in \nolinkurl{small-cut-sharp-bound.pdf}.
The unrestricted quantitative estimate appears in \nolinkurl{polynomial-cycle-error.pdf}.
These inputs are not re-proved here.

\section{A strict deficit at bounded girth}
\label{sec:path-fixed-girth}

A shortest cycle gives a sequence of terminal-path bounds.
The argument permits a disconnected complement of that cycle.

The proof uses the asymptotic cycle bound and the protected-vertices theorem.
The latter is Theorem~\ref{thm:high-density-protected-cyclic-sets}.
Thus its quantitative constants ultimately depend on additive terminal-path stability.

\begin{lemma}[Prefix deletion along a shortest cycle]
\label{lem:path-shortest-cycle-prefix}
Let $G$ be a connected loopless cubic multigraph of rank $r$ and girth $g\ge5$.
Let $C=v_0v_1\cdots v_{g-1}v_0$ be a shortest cycle.
Write $w_i$ for the unique neighbor of $v_i$ outside $C$.
Put $N=G-v_0v_{g-1}$ and $b=\beta(N)=r-1$.
For $1\le p\le g-1$, delete $v_0,\ldots,v_{p-1}$ from $N$.
Let $K_p$ be the component containing $v_{g-1}$.
Then
\[
 \beta(K_p)\le b-p.
\]
If $p\le g-2$, then $K_p$ is a normalized pole with at most $g$ ports.
If $p=g-1$, deleting the pendant vertex $v_{g-1}$ from $K_p$ gives such a pole of the same rank.
Whenever $w_{p-1}\in K_p$, it is a port of the resulting pole.
Its terminal partner is $v_{g-1}$ for $p\le g-2$, and $w_{g-1}$ for $p=g-1$.
These terminal partners are distinct from $w_{p-1}$.
\end{lemma}

\begin{proof}
The cycle $C$ is induced.
An outside vertex cannot have two neighbors on $C$.
Otherwise the shorter arc between those neighbors gives a cycle of length at most
\[
 \left\lfloor\frac g2\right\rfloor+2<g.
\]
Consequently all $w_i$ are distinct, and every component of $G-V(C)$ has minimum degree at least two.
Each such component contains a cycle.

Deleting the prefix removes $p$ cycle-path edges and $p$ outside edges.
If the remaining graph has $c_p$ components, its total rank is
\[
 b-p+c_p-1.
\]
Every component other than $K_p$ avoids the remaining cycle path.
It is therefore an entire component of $G-V(C)$ and has rank at least one.
Subtracting these $c_p-1$ ranks gives $\beta(K_p)\le b-p$.

For $p\le g-2$, the two endpoints of the remaining cycle path have degree two.
Its other vertices have degree three.
An outside vertex loses at most one incident edge, so its degree remains at least two.
The degree-two vertices are these two endpoints and some of $w_0,\ldots,w_{p-1}$.
There are at most $p+2\le g$ such vertices.
In particular, $w_{p-1}$ has degree two whenever it belongs to $K_p$.

For $p=g-1$, the remaining cycle vertex has only its outside edge.
Delete that pendant vertex and move the target to $w_{g-1}$.
The graph becomes a component of $G-V(C)$ and its rank does not change.
Its ports are some of $w_0,\ldots,w_{g-1}$.
The distinctness of the outside neighbors proves the final assertion.
\end{proof}

\begin{theorem}[Fixed-girth deficit]
\label{thm:path-fixed-girth-deficit}
Fix an integer $g\ge5$.
Let $G$ range over connected loopless cubic multigraphs of rank $r$, girth $g$, and no cyclic cut of size at most four.
Then, uniformly over these graphs,
\[
 c(G)\le 2^{r-1}-2^{r-g-1}+o_g(2^r)
 \qquad(r\to\infty).
\]
In particular, $c(G)<2^{r-1}$ for all sufficiently large $r$.
\end{theorem}

\begin{proof}
Use the notation of Lemma~\ref{lem:path-shortest-cycle-prefix}.
Let $P$ count the simple $v_0$--$v_{g-1}$ paths in $N$.
The cycle path $v_0v_1\cdots v_{g-1}$ contributes one.
Every other path first leaves this fixed path through $v_iw_i$, for some $0\le i\le g-2$.
Its remaining segment avoids $v_0,\ldots,v_i$.

Put $p=i+1$.
If $w_i\notin K_p$, the corresponding branch contributes zero.
Otherwise its contribution is exactly the terminal-path count in the pole from Lemma~\ref{lem:path-shortest-cycle-prefix}.
Write $q_p$ for that pole's rank.
Then $q_p\le b-p$.

Fix $\varepsilon>0$.
There is a threshold $B=B(g,\varepsilon)$ such that every nonzero branch of rank at least $B$ has terminal-path count at most
\[
 (1/2+\varepsilon)2^{q_p}.
\]
Indeed, apply the protected-vertices theorem for each possible number of ports, between two and $g$.
Protect all remaining vertices of $C$.
The theorem supplies two disjoint cyclic vertex sets if this bound fails.
They avoid every port and every protected vertex.
Every changed incidence in passing back to $G$ is at a port.
Both sets therefore retain their cycles and boundaries of size at most four in $G$, which is impossible.

A branch of rank less than $B$ has at most $2^B$ terminal paths.
This follows from the usual injection of terminal paths into their affine parity class.
Hence every branch, uniformly in its rank, has at most
\[
 (1/2+\varepsilon)2^{b-p}+2^B
\]
paths.
Summing over the $g-1$ branches gives
\begin{align*}
 P
 &\le 1+(1/2+\varepsilon)\sum_{p=1}^{g-1}2^{b-p}+(g-1)2^B\\
 &=1+(1/2+\varepsilon)\bigl(2^b-2^{b-g+1}\bigr)+(g-1)2^B.
\end{align*}
First let $r\to\infty$, and then let $\varepsilon\downarrow0$.
The threshold $B$ depends only on $g$ and $\varepsilon$.
Therefore
\[
 P\le 2^{b-1}-2^{b-g}+o_g(2^b).
\]

Every cycle of $G$ either avoids $v_0v_{g-1}$ or corresponds to one of these terminal paths.
Thus $c(G)=c(N)+P$.
The asymptotic cycle bound applies to the connected graph $N$ of rank $b$ and gives
\[
 c(N)\le2^{b-1}+o(2^b).
\]
Adding the last two estimates and substituting $b=r-1$ proves the theorem.
\end{proof}

\begin{corollary}[The remaining girth must tend to infinity]
\label{cor:path-unbounded-girth-reduction}
Fix an integer $g_0\ge2$.
There is a rank threshold $R(g_0)$ such that every connected loopless cubic multigraph $G$ with
\[
 \beta(G)\ge R(g_0),\qquad \operatorname{girth}(G)\le g_0,
\]
and no cyclic cut of size at most four satisfies $c(G)<2^{\beta(G)-1}$.
Consequently, any sequence with ranks tending to infinity, no such cyclic cuts, and cycle counts above the half bound has girth tending to infinity.
\end{corollary}

\begin{proof}
If a cycle has at most four vertices, deleting its vertices reduces rank by at most eight.
For rank greater than eight, its complement therefore contains a cycle.
Its boundary has at most four edges, which contradicts the cut assumption.
Thus the girth is at least five at sufficiently large rank.
Apply Theorem~\ref{thm:path-fixed-girth-deficit} to the finitely many values $5\le g\le g_0$.
\end{proof}

The argument does not require $G-V(C)$ to be connected.
It also does not require the cut around $C$ to have minimum cyclic order.
The rank bound in Lemma~\ref{lem:path-shortest-cycle-prefix} is the needed substitute for those assumptions.

\begin{corollary}[Every fixed-order cut with a bounded-rank shore]
\label{cor:path-bounded-shore-any-order}
Fix integers $k,H\ge1$.
For all sufficiently large ranks, the following statement holds.
Let $G$ be a connected loopless cubic multigraph with no cyclic cut of size at most four.
If $G$ has a cyclic cut of size at most $k$ with one shore of total rank at most $H$, then
\[
 c(G)<2^{\beta(G)-1}.
\]
Neither shore is required to be connected, and distinct cut edges may have a common endpoint.
\end{corollary}

\begin{proof}
Choose a cyclic component $A$ of the bounded-rank shore.
Write $h_A=\beta(G[A])\le H$ and $d_A=|\delta_G A|\le k$.
The cubic degree identity gives
\[
 |A|=2h_A+d_A-2\le2H+k-2.
\]
Thus $A$ contains a cycle of length at most $2H+k-2$.
Apply Corollary~\ref{cor:path-unbounded-girth-reduction} with the fixed bound $\max\{2,2H+k-2\}$.
\end{proof}

In particular, the remaining case for any fixed cut order has both shore ranks tending to infinity.
Combining this corollary with the established sharp bound for cyclic cuts of size at most four also gives
\[
 c(G)\le2^{r-1}+O_{k,H}(r2^{r/2})
\]
for every graph with such a bounded-rank cyclic shore, without the small-cut exclusion.
The finitely many ranks below the strict-deficit threshold are absorbed into the constant.

\subsection{A finite inequality and a growing-girth condition}

The branch poles already contain every changed incidence at a port.
Thus the unprotected small-cut theorem gives a slightly smaller explicit threshold.

\begin{proposition}[Finite shortest-cycle estimate]
\label{prop:path-fixed-girth-finite}
Let $G$ satisfy the hypotheses of Theorem~\ref{thm:path-fixed-girth-deficit}, with rank $r$ and girth $g\ge5$.
The value of $g$ need not be fixed here.
For $0<\varepsilon\le1/2$, put
\[
 B(g,\varepsilon)
  =(g+1)\bigl(9\lceil C_\varepsilon\rceil+8g+9\bigr),
 \qquad
 E(s)=\max\{\psi(s)-2^{s-1},0\}.
\]
Then
\begin{align*}
 c(G)\le{}&2^{r-1}-2^{r-g-1}
       +\varepsilon\bigl(2^{r-1}-2^{r-g}\bigr)
       +E(r-1)\\
       &\hspace{20mm}+1+(g-1)2^{B(g,\varepsilon)}.
\end{align*}
\end{proposition}

\begin{proof}
Each nonzero branch is a normalized pole with between two and $g$ ports.
Every incidence added in returning to $G$ is at one of its ports.
Theorem~\ref{thm:high-density-pair-small-cyclic-cut} therefore bounds each branch of rank at least $B(g,\varepsilon)$ by $(1/2+\varepsilon)2^{q_p}$.
A branch below this threshold has at most $2^{B(g,\varepsilon)}$ paths.
Use $q_p\le r-1-p$, sum the geometric series, and add
\[
 c(G-v_0v_{g-1})\le2^{r-2}+E(r-1).
\]
This gives the displayed finite inequality.
\end{proof}

For example, consider a sequence with ranks $r\to\infty$, girths $g=g(r)\ge5$, and no cyclic cut of size at most four.
Set $\varepsilon=2^{-g-2}$.
The sequence is eventually below the half bound if
\begin{equation}\label{eq:path-growing-girth-conditions}
 \frac{2^g E(r-1)}{2^r}\longrightarrow0,
 \qquad
 B(g,2^{-g-2})+g+\log_2 g-r\longrightarrow-\infty.
\end{equation}
Indeed, the $\varepsilon$ term is at most one quarter of the negative term $2^{r-g-1}$.
Both remaining error terms are negligible compared with that negative term under these conditions.

The proved endpoint cycle bound makes the first condition hold whenever
\[
 \frac{2^g(\log r)^{11/2}}{\sqrt r}\longrightarrow0.
\]
The second condition still depends on the additive stability constant $C_{2^{-g-2}}$.
No polynomial estimate for that constant is asserted here.


\section{Independent finite verification}
The standard-library script \texttt{check\_root\_shortest\_cycle\_prefix.py} checks the finite identities on seventy-two random simple cubic graphs of girth at least five.
It also checks an explicit shortest five-cycle with two cyclic exterior components.
There are 292 prefix checks in total.
For each graph, independent cycle-space enumeration verifies that deleting the selected edge removes exactly the counted terminal paths.
A separate path enumeration verifies their first-divergence partition.
The checks include the final pendant-target reduction, all port degrees, and the component-rank inequalities.

For the disconnected-complement example, the four branch counts are $812,406,0,28$.
Together with the cycle path, they give $1,247$ cycles through the selected edge.
The zero branch confirms that the target-component restriction is needed.
The graph has rank thirteen and girth five.

These finite checks do not establish the asymptotic estimate.
That conclusion follows from the proof and the analytic inputs stated above.
The second growth condition still depends on $C_\varepsilon$; no useful polynomial estimate for that constant is claimed.
\end{document}
